\section{Crossed twists, finite Fourier corrections and balanced order jets}
\label{btc:section}
The common-interval mixture in Section~\ref{mxtw:sec:main} keeps each
frequency on a single complex ray. We now return to that spectral calculus
and allow the intermediate twist to cross integers. At positive integer
total order the missing term is a finite Fourier polynomial, determined by
an affine finite part. This continuation of the beta-crossing
report~\cite{btc:report} also supplies an analytic family for the balanced
order derivatives that lead into Chapter~9. The earlier common-interval
theorem retains its stated scope.

For real nonintegral $\theta$, retain the principal-ray multiplier
\[
 \widehat K_\alpha^\theta(n)=
 \exp\!\left[-\alpha\left(\log(2\pi|n+\theta|)
                +\frac{\pi i}{2}\operatorname{sgn}(n+\theta)\right)\right].
\]
On compact order sets its coefficients and order derivatives have uniform
polynomial growth. Thus $K_\alpha^\theta$ is an entire family of periodic
distributions, and finite convolutions are defined by multiplying these
coefficients.

\subsection{The scalar finite part fixes the crossing normalization}
Let $d\ge2$, $y_j\in\mathbb R\setminus\{0\}$,
$\operatorname{Re}\alpha_j>0$ and $\sum_j\alpha_j=m\in\mathbb N_{>0}$.
On $\Sigma_{d-1}=\{t_j\ge0:\sum t_j=1\}$, with
$dt=dt_1\cdots dt_{d-1}$, put
\[
 \sigma=\sum_{y_j<0}\alpha_j,\qquad
 Q=\prod_j|y_j|^{-\alpha_j},\qquad
 d\mu(\mathbf t)=\frac{\Gamma(m)}{\prod_j\Gamma(\alpha_j)}
                         \prod_jt_j^{\alpha_j-1}\,dt.
\]
The pushforward of this finite complex measure under the \emph{affine}
coordinate $y=\sum t_jy_j$ will be denoted by $h$. It need not be positive.
If both signs occur, $h$ has a smooth density near zero. To see this,
group the negative and positive vertices, choose their internal Dirichlet
coordinates, and let $v$ be the total negative mass. Its conditional law is
\[
 \frac{\Gamma(m)}{\Gamma(\sigma)\Gamma(m-\sigma)}
                    v^{\sigma-1}(1-v)^{m-\sigma-1}\,dv.
\]
The internal weighted means satisfy $Y_-\le-c<0<c\le Y_+$ uniformly.
Since $y=vY_-+(1-v)Y_+$, the inverse
$v=(Y_+-y)/(Y_+-Y_-)$ stays away from both endpoints near zero.
All fixed $y$ derivatives of its density are bounded by integrable
internal Dirichlet weights. This proves smoothness and holomorphic
parameter dependence on compact subsets of the indicated hyperplane.
If all signs agree, the support stays away from zero.

Use the finite-part distribution
\[
 \operatorname{Pf}y^{-m}
 =\frac{(-1)^{m-1}}{(m-1)!}
              \frac{d^{m-1}}{dy^{m-1}}\operatorname{PV}\frac1y.
\]
Pairing it with $h$ means local Taylor subtraction near zero and an
ordinary measure integral elsewhere. The cutoff variable is exactly
$y$; a nonlinear substitution without its local correction is not allowed.
The boundary identity is
\begin{equation}\label{btc:boundary}
 (y-i0)^{-m}=\operatorname{Pf}y^{-m}
       +\frac{i\pi(-1)^{m-1}}{(m-1)!}\delta^{(m-1)}(y).
\end{equation}
Indeed $\varepsilon/(y^2+\varepsilon^2)\to\pi\delta$ and
$y/(y^2+\varepsilon^2)\to\operatorname{PV}(1/y)$; differentiation
gives every $m$. Consequently its pairing adds
$i\pi h^{(m-1)}(0)/(m-1)!$.

\begin{theorem}[Dirichlet crossing finite part]\label{btc:scalar}
In this normalization,
\begin{align}
 \operatorname{FP}\int h(y)y^{-m}\,dy&=\cos(\pi\sigma)Q,
                                      \label{btc:scalar-fp}\\
 h^{(m-1)}(0)&=\frac{(m-1)!}{\pi}\sin(\pi\sigma)Q.
                                      \label{btc:scalar-density}
\end{align}
In a one-sign case the density derivative is zero, with $\sigma=0$ or $m$.
\end{theorem}
\begin{proof}
For $\operatorname{Re}z_j>0$, multiply the Gamma integrals and substitute
$u_j=rt_j$. Absolute convergence and the Jacobian $r^{d-1}$ give
\[
 \prod_jz_j^{-\alpha_j}=
       \int_{\Sigma_{d-1}}\left(\sum_jt_jz_j\right)^{-m}d\mu.
\]
Put $z_j=\varepsilon+i y_j$ and cancel $i^{-m}$. As
$\varepsilon\downarrow0$ the left side becomes $e^{i\pi\sigma}Q$;
each negative vertex contributes $e^{i\pi\alpha_j}$. Smoothness at zero,
\eqref{btc:boundary}, and dominated convergence away from zero justify the
limit on the right. Taking $z_j=\varepsilon-i y_j$ similarly yields
the other boundary value $e^{-i\pi\sigma}Q$. Averaging gives
\eqref{btc:scalar-fp}; subtracting gives \eqref{btc:scalar-density}.
The parameters may be complex: the two limits are not being treated as
complex conjugates.
\end{proof}

For $A,B>0$, the two-vertex case is the explicit identity
\begin{equation}\label{btc:two-scalar}
 \frac{\Gamma(m)}{\Gamma(\alpha)\Gamma(\beta)}
 \operatorname{FP}\int_0^1
 \frac{t^{\alpha-1}(1-t)^{\beta-1}}{[B-(A+B)t]^m}\,dt
       =\frac{\cos(\pi\alpha)}{A^\alpha B^\beta},
 \qquad \alpha+\beta=m.
\end{equation}
Its crossing is $t=B/(A+B)$ and the omitted boundary term is
$i\sin(\pi\alpha)/(A^\alpha B^\beta)$.

\subsection{Only the crossed Fourier modes need correction}
For real nonintegral twists $\theta_j$, define
\[
 \mathcal C=\{n\in\mathbb Z:
       \min_j(n+\theta_j)<0<\max_j(n+\theta_j)\},\qquad
 \sigma_n=\sum_{n+\theta_j<0}\alpha_j,\qquad
 Q_n=\prod_j|n+\theta_j|^{-\alpha_j}.
\]
Let $\mathcal M_m$ be the simplex mixture whose $n$th coefficient is
\[
 (2\pi i)^{-m}\operatorname{FP}
          \int_{\Sigma_{d-1}}\left(n+\sum_jt_j\theta_j\right)^{-m}d\mu.
\]
This defines the finite-part interpretation of the intermediate kernel
$K_m^{\sum t_j\theta_j}$. Only finitely many coefficients cross zero.
\begin{theorem}[Integer-total mixture across arbitrary twists]
\label{btc:periodic}
For $\operatorname{Re}\alpha_j>0$ and $\sum\alpha_j=m\in\mathbb N_{>0}$,
\begin{equation}\label{btc:mixture}
 K_{\alpha_1}^{\theta_1}*\cdots*K_{\alpha_d}^{\theta_d}
 =\mathcal M_m+\frac{i}{(2\pi i)^m}
      \sum_{n\in\mathcal C}\sin(\pi\sigma_n)Q_ne^{2\pi i n x}.
\end{equation}
Both terms are holomorphic on the indicated parameter hyperplane.
The correction vanishes for a common open unit interval, and also for
every positive-integer order tuple even when crossings occur.
\end{theorem}
\begin{proof}
The product multiplier is $(2\pi i)^{-m}e^{i\pi\sigma_n}Q_n$,
while Theorem~\ref{btc:scalar} gives the mixture coefficient
$(2\pi i)^{-m}\cos(\pi\sigma_n)Q_n$. Their difference is
\eqref{btc:mixture}. Off $\mathcal C$, $\sigma_n$ is $0$ or $m$.
For large $|n|$, all shifted frequencies have one sign and are comparable
to $|n|$. Coefficients are $O(|n|^{-m})$ uniformly on compact parameter
sets; order derivatives add at most powers of $\log|n|$.
Pairing with smooth test functions justifies all series and derivatives,
and Fourier uniqueness proves equality. At $m\ge2$ the undifferentiated
series is absolutely summable; distributions also cover $m=1$.
\end{proof}
The correction is a finite trigonometric polynomial in $x$. The delta
derivative in \eqref{btc:boundary} belongs to the intermediate affine
twist coordinate, not to the spatial variable $x$.

For $a<b$ nonintegral and $\alpha+\beta=m$, the correction is
\[
 \frac{i\sin(\pi\alpha)}{(2\pi i)^m}
 \sum_{\substack{k\in\mathbb Z\\a<k<b}}
       \frac{e^{-2\pi i kx}}{(k-a)^\alpha(b-k)^\beta}.
\]
Repeated twists combine by adding their orders, directly from either the
Dirichlet grouping or the Fourier multipliers.
At $a=-1/2$, $b=1/2$, $\alpha=\beta=1/2$, one obtains
\begin{equation}\label{btc:counterexample}
 K_{1/2}^{-1/2}*K_{1/2}^{1/2}
 =\frac1\pi\operatorname{PV}\int_0^1
          \frac{K_1^{t-1/2}}{\sqrt{t(1-t)}}\,dt+\frac1\pi.
\end{equation}
The convolution has zero-mode coefficient $1/\pi$, while the uncorrected
mixture has zero. This is an exact obstruction to extending the old
mixture without the common-interval hypothesis.

\subsection{Balanced derivatives and the exact chamber kernel}
Fix positive integers $r_j$ with sum $m$, a direction $\mathbf v$ with
$\sum v_j=0$, and $\alpha_j=r_j+uv_j$. For a crossed mode put
\[
 R_n=\sum_{n+\theta_j<0}r_j,\quad
 V_n=\sum_{n+\theta_j<0}v_j,\quad
 L_n=-\sum_jv_j\log|n+\theta_j|,\quad
 Q_{n,0}=\prod_j|n+\theta_j|^{-r_j}.
\]
\begin{theorem}[All balanced correction jets]\label{btc:jets}
The $q$th $u$ derivative of the correction at zero is
\begin{equation}\label{btc:jet}
 C_q(x)=\frac1{2(2\pi i)^m}\sum_{n\in\mathcal C}
 (-1)^{R_n}Q_{n,0}e^{2\pi i n x}
 \bigl[(L_n+i\pi V_n)^q-(L_n-i\pi V_n)^q\bigr].
\end{equation}
For several balanced directions, the bracket becomes the difference of
the corresponding products of powers. If the twists occupy $c$ distinct
open unit intervals with index groups $I_1,\ldots,I_c$, the map
$\mathbf v\mapsto C_1$ has rank $c-1$ on the balanced tangent space.
Its kernel has dimension $d-c$ and is exactly
\begin{equation}\label{btc:chamber-kernel}
 \left\{\mathbf v:\sum_{j\in I_h}v_j=0\text{ for every }h\right\}.
\end{equation}
These are precisely the directions in which every correction jet
vanishes and the full correction is identically zero near $u=0$.
\end{theorem}
\begin{proof}
At a crossed mode,
\[
 i\sin(\pi(R_n+uV_n))Q_{n,0}e^{uL_n}
 =\frac{(-1)^{R_n}Q_{n,0}}2
    \bigl[e^{u(L_n+i\pi V_n)}-e^{u(L_n-i\pi V_n)}\bigr].
\]
Differentiate this finite sum within the analytic family already proved.
Distinct frequencies are linearly independent and $Q_{n,0}\ne0$, so
$C_1=0$ exactly when every $V_n=0$. At $n=-k$ the negative group is
$\{j:\theta_j<k\}$: as $k$ runs through the crossed integers these
are precisely the proper prefixes of the occupied chambers, with possible
repetitions across gaps. Their zero sums, together with balance, are
equivalent to \eqref{btc:chamber-kernel}. Chamber sums map onto the
$(c-1)$-dimensional zero-sum space, proving the rank. For a kernel
direction each sine factor remains identically zero. The converse
already follows from the first derivative.
\end{proof}
For the orders $(1+u,1-u)$ at twists $(-1/2,1/2)$, the zero-mode
correction is $i\sin(\pi u)/\pi^2$. Its value at zero is zero but
its first derivative is $i/\pi$. An identity restricted to integer
orders therefore cannot be differentiated to omit this correction.

For completeness, the simplex weight itself has an explicit derivative
compiler. At the integer tuple put
$B_0=\Gamma(m)\prod t_j^{r_j-1}/\prod\Gamma(r_j)$. Its $q$th balanced
derivative is $B_0Y_q(\ell_1,\ldots,\ell_q)$, where $Y_q$ is the complete
exponential Bell polynomial and
\begin{align*}
 \ell_1&=\sum_jv_j(\log t_j-H_{r_j-1}),\\
 \ell_k&=(-1)^{k+1}(k-1)!\sum_jv_j^k
                  (\zeta(k)-H_{r_j-1}^{(k)}),\qquad k\ge2.
\end{align*}
Differentiate the logarithm of the density, use the integer polygamma
series, and exponentiate; the Euler constant cancels by balance.
These are genuine order derivatives of an analytic family, with the
same affine finite-part normalization. The chamber rank is a spectral
rank and makes no arithmetic-independence claim about point values.

\subsection{Ordinary centered logarithmic moments reduce to even zeta values}
The same affine identity yields ordinary integrals, rather than only
distributional corrections. For $A,B>0$ put
\[
 U(t)=\log\frac{At}{B(1-t)},\qquad D(t)=B-(A+B)t.
\]
Both vanish simply at $t_0=B/(A+B)$; the centered logarithm is essential.
For positive integers $r,s$, write $m=r+s$ and
\[
 P_{r,s}(u)=\prod_{j=1}^{r-1}(j+u)\prod_{j=1}^{s-1}(j-u).
\]
\begin{theorem}[All centered logarithmic moments]\label{btc:moments}
For $-r<\operatorname{Re}u<s$, the affine finite-part generator is
\begin{equation}\label{btc:moment-generator}
 \operatorname{FP}\int_0^1
 \frac{t^{r-1}(1-t)^{s-1}e^{uU(t)}}{D(t)^m}\,dt
 =\frac{(-1)^r\pi u\cot(\pi u)P_{r,s}(u)}
              {(m-1)!A^rB^s}.
\end{equation}
The singularities inside this strip on the right are removable. For
every integer $q\ge m$, the following integral is ordinary and absolutely
convergent, and has the finite even-zeta evaluation
\begin{equation}\label{btc:ordinary-moments}
 \int_0^1\frac{t^{r-1}(1-t)^{s-1}U(t)^q}{D(t)^m}\,dt
 =\frac{2(-1)^{r+1}q!}{(m-1)!A^rB^s}
       \sum_{k=1}^{\lfloor q/2\rfloor}
          \zeta(2k)[u^{q-2k}]P_{r,s}(u).
\end{equation}
In particular, for $k\ge1$,
\begin{equation}\label{btc:even-logarithms}
 \int_0^1\frac{\log^{2k}(t/(1-t))}{(1-2t)^2}\,dt
                       =2(2k)!\zeta(2k).
\end{equation}
\end{theorem}
\begin{proof}
Apply \eqref{btc:two-scalar} with $\alpha=r+u$ and $\beta=s-u$.
The centering factor $(A/B)^u$ cancels the varying powers on the right.
Gamma recurrence and reflection then give
\[
 \frac{\Gamma(r+u)\Gamma(s-u)\cos(\pi(r+u))}
         {(m-1)!A^rB^s}
 =\frac{(-1)^r\pi u\cot(\pi u)P_{r,s}(u)}{(m-1)!A^rB^s}.
\]
The original Gamma expression is holomorphic in the stated strip, so the
apparent cotangent poles cancel. Taylor subtraction near $t_0$, and
endpoint integrability uniformly on compact substrips, justify the order
derivatives. At $q\ge m$, $U(t)^q/D(t)^m=O(|t-t_0|^{q-m})$;
the derivative is an ordinary integral with integrable endpoint logarithms.
The partial fractions of cotangent give the normally convergent expansion
\[
 \pi u\cot(\pi u)=1-2\sum_{k\ge1}\zeta(2k)u^{2k},\qquad |u|<1.
\]
Since $\deg P_{r,s}=m-2<q$, its initial polynomial contributes no
$u^q$ coefficient. Coefficient extraction proves
\eqref{btc:ordinary-moments}; $r=s=A=B=1$ gives
\eqref{btc:even-logarithms}.
\end{proof}
There is also an ordinary Mellin route, independently useful for experiments.
The substitution $v=|U(t)|$ on the two halves of the interval gives
\[
 \frac1{A^rB^s}\int_0^\infty
 \frac{v^q\bigl((-1)^m e^{-sv}+(-1)^q e^{-rv}\bigr)}
                         {(1-e^{-v})^m}\,dv.
\]
At $q\ge m$, both individual terms are integrable. Expanding
$(1-e^{-v})^{-m}$ gives a binomially weighted convergent Hurwitz sum.
Its polynomial multiplicities reduce it to ordinary zeta values, and
the two halves cancel all odd-zeta terms. The independent verifier checks
this finite coefficient reduction as well as direct Mellin quadrature.
The later logarithmic Mellin finite parts and Stieltjes primitives in the
delivery retain their separate audit scope.
